\subsection{An exact junction calculation for strict form truncations}

\paragraph{IC01 (Proved).}
Let $f:P\to A$ be a morphism of nonnegative cochain complexes of
simplicial $k$-modules. Fix an integer $d\geq1$. Write
$P_{<d}$ and $A_{<d}$ for the brutal quotient complexes in form
degrees strictly below $d$, and form the pullback
\[
 E=P_{<d}\mathop{\times}_{A_{<d}} A
\]
in cochain complexes of simplicial $k$-modules. Then the underlying
terms are $E^q=P^q$ for $q<d$ and $E^q=A^q$ for $q\geq d$.
The junction differential is
$d_E^{d-1}=d_A^{d-1}f^{d-1}:P^{d-1}\to A^d$.
Consequently, simplexwise and also as simplicial modules,
\[
 H^{d-1}(E)=
 \frac{\ker(d_A^{d-1}f^{d-1})}
      {\operatorname{im}(d_P^{d-2})},\qquad
 H^d(E)=
 \frac{\ker(d_A^d)}
      {f^d(\operatorname{im}(d_P^{d-1}))}.
\]
Here $P^{-1}=0$. In the category of strict PD differential
algebras the same pullback exists for the quotient by the ideal of
forms of degree at least $d$, and its underlying cochain complex
has exactly these formulas.

\begin{proof}
In each form degree below $d$, the map $A\to A_{<d}$ is the
identity. A pair in the pullback therefore has the unique form
$(p,f(p))$, giving the stated identification with $P^q$.
In degree at least $d$, both truncated terms vanish, so the
pullback term is $A^q$. The differential of $(p,f(p))$ at the
junction has first component zero and second component
$d_Af(p)=f(d_Pp)$. Below and above the junction it is the
original differential. Taking kernels and images gives the
displayed formulas. The ideals of forms of degree at least $d$
are differential PD ideals: differential raises form degree,
products add it, and positive divided powers multiply it.
Their quotient and pullback therefore have the stated algebraic
structures, with the same underlying complex.
\end{proof}

\paragraph{IC02 (Proved; explicit failure of exactness at the junction).}
For every nonzero commutative ring $k$ and every $d\geq1$, the
preceding pullback can have nonzero cohomology in degree $d-1$
even when $P$ and $A$ are contractible complexes and $f$ is a
quasi-isomorphism. Let $P$ be the identity disk
$k\xrightarrow{1}k$ in degrees $d-1,d$, let $A=0$, and let
$f=0$. Then $E=k[d-1]$ and $H^{d-1}(E)=k$.

\begin{proof}
The disk has the evident contracting homotopy, and the zero
complex is contractible. Both have zero cohomology, so the unique
map is a quasi-isomorphism. Its strict splice is the brutal
truncation of the disk, which consists of $k$ in degree $d-1$
and zero elsewhere. This proves the asserted cohomology.
\end{proof}

This counterexample is a statement about the underlying complex
construction. It makes no claim to refute the admissible
common-source conjecture. The computation prevents us from using
partial quasi-isomorphism alone to declare this strict splice an
admissible completion.

\paragraph{IC03 (Proved; strict splicing fails for admissible integral theories).}
Let $S=\Gamma_{\mathbb Z}(T)$ and let $F$ be the standard
PD de Rham theory. Choose integers $d\geq2$ and $r\geq d$.
Let $D$ be the free identity cochain disk with basis $e,u$,
degrees $d-1,d$, weights $r,r$, and differential $de=u$.
Put
\[
 P=F\oplus(F\otimes_{\mathbb Z}D),\qquad f:P\longrightarrow F.
\]
The second summand is a square-zero differential ideal; its
higher divided powers vanish. The map $f$ is projection.
Then $P$ and $F$ are admissible integral theories with free
homogeneous terms and support $q\leq w$. The projection is a
simplexwise vertical cochain homotopy equivalence. Nevertheless
\[
 E=P_{<d}\mathop{\times}_{F_{<d}}F
\]
fails eventual vertical exactness. In every weight $w\geq r$
its simplicial degree-zero cohomology in form degree $d-1$
contains a direct summand $S^{w-r}e\cong\mathbb Z$.

\begin{proof}
The square-zero product uses the graded $F$-module structure of
$F\otimes D$. For an even element $a+m$ in the PD ideal, with
$a\in F$ and $m\in F\otimes D$, define
\[
 \gamma_j(a+m)=\gamma_j(a)+\gamma_{j-1}(a)m
 \quad(j\geq2).
\]
All products of two module elements vanish. Substitution in the
PD addition, scalar, product, and composition formulas therefore
reduces their terms of module degree zero to the PD formulas in
$F$ and their terms of module degree one to the corresponding
linear terms of the same identities. The identities also follow
by adjoining formal variables for the finitely many entries and
checking the usual PD polynomial identities modulo the square
of the module ideal. Odd squares and higher divided powers of
products of odd elements vanish. The differential compatibility
reduces in the same way to the Leibniz rule for the module.
Thus this is a strict PD differential algebra over $S$.

The tensor disk is contractible. Explicitly its degree-minus-one
homotopy is $h(a u)=(-1)^{|a|}a e$ and $h(a e)=0$ for homogeneous
$a\in F$; direct substitution gives $dh+hd=1$. It is weight
preserving and commutes with simplicial operators. The inclusion
$F\to P$ and the projection $f$ consequently give a vertical
homotopy equivalence in every simplex and weight.

For a fixed form degree $q$, the extra simplicial row is the
direct sum of $F^{w-r,q-d+1}e$ and $F^{w-r,q-d}u$, omitting
negative form degrees. A standard row $F^{v,j}$ is horizontally
acyclic for $v>j$: in
$\Gamma^{v-j}(\Phi)\otimes\Lambda^j(\Phi/\mathbb ZT)$ the first
factor has a positive divided-power exponent and has the
additive simplicial contraction induced from the standard
contractible marked module $\Phi$; its terms are free.
Tensoring that contraction with the other factor preserves a
contraction. Hence for every fixed range $q\leq b$ the bound
$w\geq r+b+1$ suffices for both extra rows, and also for the
original $F$ row. This proves uniform eventual row acyclicity.
The vertical homotopy equivalence supplies the other
admissibility condition. Freeness and support follow directly
from $r\geq d$ and the corresponding properties of $F$.

Below form degree $d$, the ideal $F\otimes D$ has only its
term $F^{*,0}e$ in degree $d-1$. The strict splice discards
the entire ideal in degrees at least $d$, including its outgoing
differential. Its complex is therefore the direct sum of $F$
and this module placed in degree $d-1$ with zero differential.
For the zero-simplex, $F_0=S$, which gives the displayed
nonzero cohomology in every weight $w\geq r$. Thus no
eventual vertical bound exists.
\end{proof}
